\subsection{Convergent series}

4.1.1. Let \(f = \sum_{\alpha} f_{\alpha}\) be a formal series belonging to \(\hat{\mathbf{P}}(\mathbf{E}_1, \ldots, \mathbf{E}_n; \mathbf{F})\) , (cf. Appendix). If \(\gamma\) is a continuous semi-norm on \(\mathbf{F}\) , and \(\mathbf{R} = (\mathbf{R}_i)\) is a system of \(n\) real numbers \(>0\) , we set

\[
\| f \| _ {\gamma , \mathbf {R}} = \sup _ {\alpha} \| f _ {\alpha} \| _ {\gamma} \mathbf {R} ^ {\alpha}.
\]

The definitions and results of no. 3.1.1 (second paragraph) and of no. 3.1.2 apply without change; in particular, one defines the spaces

\[
\mathcal {H} _ {\mathrm{R}} (\mathrm{E} _ {1}, \dots , \mathrm{E} _ {n}; \mathrm{F}) \quad \text { and } \quad \mathcal {H} (\mathrm{E} _ {1}, \dots , \mathrm{E} _ {n}; \mathrm{F}).
\]

4.1.2. The canonical isomorphism \(j\) of \(\hat{\mathbf{P}}(\mathbf{E};\mathbf{F})\) onto \(\hat{\mathbf{P}}(\mathbf{E}_{1},\ldots,\mathbf{E}_{n};\mathbf{F})\) gives, by restriction, an isomorphism of topological vector spaces from \(\mathcal{H}_{\mathbb{R}}(\mathbf{E};\mathbf{F})\) onto \(\mathcal{H}_{(\mathbb{R},\ldots,\mathbb{R})}(\mathbf{E}_{1},\ldots,\mathbf{E}_{n};\mathbf{F})\) for all \(\mathbf{R}\in\mathbb{R}_{+}^{*}\) ; it also gives an isomorphism of \(\mathcal{H}(\mathbf{E};\mathbf{F})\) onto \(\mathcal{H}(\mathbf{E}_{1},\ldots,\mathbf{E}_{n};\mathbf{F})\) . More precisely, if \(f=\sum_{m}f_{m}\in\hat{\mathbf{P}}(\mathbf{E};\mathbf{F})\) and if \(j(f)=\sum_{\alpha}f_{\alpha}\) , one has, for every continuous semi-norm \(\gamma\) on \(\mathbf{F}\) :

\[
\begin{array}{l} \| f _ {m} \| _ {\gamma} = \sup _ {| \alpha | = m} \| f _ {\alpha} \| _ {\gamma} \\ \| f \| _ {\gamma , \mathbf {R}} = \| j (f) \| _ {\gamma , (\mathbf {R}, \dots , \mathbf {R})}. \end{array}
\]

4.1.3. Let \(f = \sum_{\alpha} f_{\alpha}\) be an element of \(\mathcal{H}(E_1, \ldots, E_n; F)\); let \(I(f)\) be the set of \(R \in (\mathbf{R}_+^*)^n\) such that, for every continuous semi-norm \(\gamma\) on F, the product \(\|f_{\alpha}\|_{\gamma} R^{\alpha}\) tends to zero when \(|\alpha|\) tends to infinity. The set \(I(f)\) is nonempty; it is called the strict convergence indicator of \(f\). The set \(\Omega(f)\) of points

\[
(\log \mathbf {R} _ {1}, \dots , \log \mathbf {R} _ {n}) \quad \text { for } \quad \mathbf {R} \in I (f)
\]

is a convex subset of \(\mathbf{R}^{n}\) .

When \(n=1\) , the set \(\mathbf{I}(f)\) is an interval of \(\mathbf{R}\) , open on the left and open or closed on the right; its upper bound (finite or \(+\infty\) ) is denoted \(\rho(f)\) and is called the strict radius of convergence of \(f\) .

With the notations of 4.1.2, one has \(\mathbf{R} \in \mathbf{I}(f)\) if and only if

\[
(\mathbf {R}, \dots , \mathbf {R}) \in \mathrm{I} (j (f)).
\]

The set of points \(x = (x_i)\) such that there exists \(\mathbf{R} = (\mathbf{R}_i) \in I(f)\) with \(\| x_i \| \leqslant \mathbf{R}_i\) for \(1 \leqslant i \leqslant n\) is called the strict domain of convergence of \(f\) and is denoted \(C(f)\) . It is an open subset of \(E\), the union of the polyballs

\[
\mathrm{B} (\mathbf {R}) = \{x \in \mathrm{E} | \| x _ {i} \| \leqslant \mathrm{R} _ {i} \text {   for   } 1 \leqslant i \leqslant n \},
\]

for \(\mathbf{R}\in\mathbf{I}(f)\) .

4.1.4. The results of 3.1.7 and 3.1.8 remain valid, replacing everywhere therein \(\tilde{\mathbf{C}}(f)\) by \(\mathbf{C}(f)\) and \(\tilde{\mathcal{H}}_{\mathbb{R}}\) by \(\mathcal{H}_{\mathbb{R}}\) .

4.1.5. Let \(\mathbf{F}_1, \ldots, \mathbf{F}_m\) be complete normed spaces and suppose \(\mathbf{F}\) quasi-complete. Let \(\mathbf{f} = (f_i)_{1 \leqslant i \leqslant m}\) , with \(f_i \in \mathcal{H}(\mathbf{E}_1, \ldots, \mathbf{E}_n; \mathbf{F}_i)\) and let \(g \in \mathcal{H}(\mathbf{F}_1, \ldots, \mathbf{F}_m; \mathbf{F})\) , such that the point \((f_i(0))_{1 \leqslant i \leqslant m}\) of \(\mathbf{E}\) belongs to the strict domain of convergence of \(g\) . Then, for every \(\alpha \in \mathbb{N}^m\) , the formal series \(g_\alpha \circ \mathbf{f}\) belongs to \(\mathcal{H}(\mathbf{E}_1, \ldots, \mathbf{E}_n; \mathbf{F})\) and the family of \(g_\alpha \circ \mathbf{f}\) is summable in \(\mathcal{H}(\mathbf{E}_1, \ldots, \mathbf{E}_n; \mathbf{F})\) (hence a fortiori in \(\hat{\mathbf{P}}(\mathbf{E}_1, \ldots, \mathbf{E}_n; \mathbf{F})\) ). Its sum will be denoted \(g \circ \mathbf{f}\) .

More precisely, there exists \(\mathbb{R} \in \bigcap_{i} \mathrm{I}(f_i)\) and \(\mathbb{R}' \in \mathrm{I}(g)\) such that \(\sup_{|\alpha| > 0} \| f_{i,\alpha} \| \mathbb{R}^{\alpha} < \mathbb{R}_i'\) (for \(1 \leqslant i \leqslant m\) ). Under these conditions, the formal series \(g_{\alpha} \circ \mathbf{f}\) belongs to \(\mathcal{H}_{\mathbb{R}}(\mathbf{E}_1, \ldots, \mathbf{E}_n; \mathbf{F})\) and the family of \(g_{\alpha} \circ \mathbf{f}\) is summable in \(\mathcal{H}_{\mathbb{R}}(\mathbf{E}_1, \ldots, \mathbf{E}_n; \mathbf{F})\) . Finally, if \(x \in \mathbf{B}(\mathbb{R})\) , then \(\mathbf{f}(x) = (f_i(x))\) belongs to \(\mathbf{C}(g)\) and we have:

\[
g (\mathbf {f} (x)) = (g \circ \mathbf {f}) (x).
\]

Suppose further that, for each \(i\) , there exists a family \((e_{j}^{i})\) of elements of \(\mathbf{E}_{i}\) such that every element \(x\) of \(\mathbf{E}_{i}\) is the sum of a summable family \((\lambda_{j}e_{j}^{i})\) (with \(\lambda_{j}\in\mathbf{K}\) ) such that \(\|x\|=\sup_{j}|\lambda_{j}|\) . One can then, in the preceding paragraph, replace the condition \(\sup_{|\alpha|>0}\|f_{i,\alpha}\|\mathbf{R}^{\alpha}<\mathbf{R}_{i}^{\prime}\) by the condition \(\|f_{i}\|_{\mathbf{R}}\leqslant\mathbf{R}_{i}^{\prime}\) .

4.1.6. Suppose that \(E_{i} = K\) for \(1 \leqslant i \leqslant n\) . The space \(\hat{\mathbf{P}}(\mathbf{K}^{n}; \mathbf{F})\) is then identified with the space of formal power series in n indeterminates \(X_{1}, \ldots, X_{n}\) with coefficients in F, and an element f of \(\hat{\mathbf{P}}(\mathbf{K}^{n}; \mathbf{F})\) is written:

\[
f = \sum_ {\alpha} X ^ {\alpha} c _ {\alpha} \quad \text { with } c _ {\alpha} \in F.
\]

\[
\| f \| _ {\gamma , \mathbf {R}} = \sup \| c _ {\alpha} \| _ {\gamma}. \mathbf {R} ^ {\alpha}.
\]

If \(\mathbf{R} \in (\mathbf{R}_{+}^{*})^{n}\) and if \(\gamma\) is a continuous semi-norm on \(F\) , we have:

The last paragraph of 3.1.10 remains valid, as does 3.1.11.

4.1.7. Suppose that K is a closed subfield of a complete valued field (necessarily ultrametric) L. For \(y \in E_{i} \otimes_{K} L\) , set:

\[
\| y \| = \inf (\sup _ {k} | a _ {k} |. \| x _ {k} \|)
\]

the infimum being taken over all finite families of pairs \((x_{k}, a_{k}) \in E_{i} \times L\) such that \(y = \sum_{k} x_{k} \otimes a_{k}\) . We thus obtain a norm on the L-vector space \(E_{i} \otimes_{K} L\) , which induces on \(E_{i}\) the given norm. We denote by \(E_{i}^{L}\) the completion of \(E_{i} \otimes_{K} L\) for this norm.

Now let F be a separated complete polynormed L-vector space. For every K-polynomial-continuous map \(f_{\alpha}\), homogeneous of multidegree \(\alpha\), on \(E_1 \times \cdots \times E_n\), with values in the K-vector space polynormed \(F_K\) underlying F, there exists one and only one L-polynomial-continuous map \(\tilde{f}_{\alpha}\), homogeneous of the same multidegree, on \(E_1^L \times \cdots \times E_n^L\), with values in F, which extends \(f_{\alpha}\). For every continuous semi-norm on the L-vector space F, we have

\[
\| \tilde {f} _ {\alpha} \| _ {\gamma} = \| f _ {\alpha} \| _ {\gamma}.
\]

If \(f = \sum_{\alpha} f_{\alpha} \in \mathcal{H}(E_1, \ldots, E_n; F_K)\) , then \(\tilde{f} = \sum_{\alpha} \tilde{f}_{\alpha} \in \mathcal{H}(E_1^L, \ldots, E_n^L; F)\) . The series \(f\) and \(\tilde{f}\) have the same strict convergence indicator (and the same strict radius of convergence when \(n = 1\) ).

Conversely, let L be a non-discrete closed subfield of K and let \(E_{i}^{0}\) and \(F^{0}\) be the spaces over L obtained by restriction of scalars from the \(E_{i}\) and from F. If \(f = \sum f_{\alpha} \in \mathcal{H}(E_{1}, \ldots, E_{n}; F)\) , then \(f_{\alpha} \in \mathrm{P}_{\alpha}(E_{1}^{0}, \ldots, E_{n}^{0}; F^{0})\) ; if we set \(f^{0} = \sum_{\alpha} f_{\alpha} \in \hat{\mathrm{P}}(E_{1}^{0}, \ldots, E_{n}^{0}; F^{0})\) , then \(f^{0} \in \mathcal{H}(E_{1}^{0}, \ldots, E_{n}^{0}; F^{0})\) . One has \(C(f) \subset C(f^{0})\) and \(f(x) = f^{0}(x)\) for all \(x \in C(f)\) .

